\documentclass[12pt]{article}
\usepackage{tabularx}

% Use algorithm and algpseudocode packages so \State and related commands are available

\newcommand{\duedate}{---}
\newcommand{\assignment}{MARL (Lec6) Summary} % Change to "Problem Set X"

% Change the following to your name and UNI.
\newcommand{\name}{Aksel Kretsinger-Walters, adk2164}
\newcommand{\email}{adk2164@columbia.edu}

% NOTE: Defining collaborators is optional; to not list collaborators, comment out the
% line below. Maximum of two collaborators per problem set.
%\newcommand{\collaborators}{Vig Vigerton (\texttt{UNI}), Alice Bob (\texttt{UNI})}
% No collaborators on PS 0

\makeatletter
\def\input@path{{../}{../../}{../../../}} % add as many parents as you need
\makeatother
\input{pset_template_RL.tex} %% DO NOT CHANGE THIS LINE

\begin{document}
\psetheader %% DO NOT CHANGE THIS LINE

\section*{Lecture Summary: Multi-Agent Reinforcement Learning (MARL)}

\subsection*{1. Motivation}
In classical RL, the environment is stationary and external.
In \emph{multi-agent} settings, other agents \emph{are} part of the environment:
they learn, adapt, and respond.
Thus, single-agent RL assumptions break down, and game-theoretic ideas become essential.

\section*{2. Game Theory Basics}

\subsection*{One-Shot Games}
A normal-form game consists of:
\begin{itemize}
		\item players $1,\dots,n$,
		\item action sets $A_i$,
		\item utility functions $u_i: A \to \mathbb{R}$ where $A = A_1 \times \cdots \times A_n$.
\end{itemize}

\paragraph{Best response.}
Given actions of others $a_{-i}$,
\[
		BR_i(a_{-i}) = \arg\max_{a_i \in A_i} u_i(a_i, a_{-i}).
\]

\paragraph{Dominant strategy.}
Action $a_i$ \emph{weakly dominates} $a'_i$ if
\[
		u_i(a_i,a_{-i}) \ge u_i(a'_i,a_{-i}) \quad \forall a_{-i}.
\]
If such an action exists, it is optimal regardless of the opponent’s behavior.

\subsection*{Solution Concepts}

\paragraph{Dominant Strategy Equilibrium (DSE).}
Each player's action is a dominant strategy.

\paragraph{Nash Equilibrium (NE).}
An action profile $a=(a_1,\dots,a_n)$ is a NE if every player is playing a best
response:
\[
		u_i(a_i,a_{-i}) \ge \max_{a'_i} u_i(a'_i,a_{-i}).
\]

\paragraph{Mixed-Strategy NE.}
Each player uses a distribution $p_i \in \Delta(A_i)$ over actions.
A profile $p=(p_1,\dots,p_n)$ is NE if:
\[
		u_i(p_i,p_{-i}) \ge u_i(a_i,p_{-i}) \quad \forall a_i.
\]

\paragraph{Nash’s Theorem.}
Every finite game has a mixed-strategy Nash equilibrium.

\subsection*{Correlated and Coarse-Correlated Equilibria}
Mixed NE require independent sampling across players.
\emph{Correlated} equilibria allow a joint distribution $D$ over actions.

\paragraph{Correlated equilibrium (CE).}
No player benefits from deviating given the recommended action $a_i$:
\[
		\mathbb{E}_{a\sim D}[u_i(a_i,a_{-i})]
		\ge
		\mathbb{E}_{a\sim D}[u_i(a'_i,a_{-i}) | a_i].
\]

\paragraph{Coarse CE (CCE).}
Relaxed condition: no player benefits by switching to a fixed alternative action.

\paragraph{Hierarchy.}
\[
		\text{DSE} \subset \text{PSNE} \subset \text{MSNE} \subset \text{CE} \subset \text{CCE}.
\]

\section*{3. Zero-Sum Games and Minimax}

For two-player zero-sum games ($u_1 = -u_2$),
the utility can be represented by a matrix $U$.

\subsection*{Minimax Theorem (Von Neumann)}
\[
		\min_q \max_p \; p^\top U q
		=
		\max_p \min_q \; p^\top U q.
\]
The value of the game is $v^*$, and any minimax pair $(p^*,q^*)$ is a Nash equilibrium.

This generalizes to continuous-action games when $u(x,y)$ is convex in $x$ and concave in $y$.

\section*{4. Markov (Stochastic) Games}

\paragraph{Definition.}
A two-player zero-sum Markov game consists of:
\[
		(S, A_1, A_2, P, R),
\]
where both players act simultaneously at each step.

\paragraph{Policies.}
Each player has a Markov policy $\pi_i: S \to \Delta(A_i)$.

\paragraph{Value of policy pair.}
For discounted infinite horizon,
\[
		V^{\pi_1,\pi_2}(s)
		=
		\mathbb{E}\left[\sum_{t=1}^\infty \gamma^{t-1} R(s_t,a_{1,t},a_{2,t}) \mid s_1=s\right].
\]

\subsection*{Minimax for Markov Games (Shapley)}
There exists a value function $V^*$ such that
\[
		V^*(s) = \max_{\pi_1} \min_{\pi_2} V^{\pi_1,\pi_2}(s)
		= \min_{\pi_2} \max_{\pi_1} V^{\pi_1,\pi_2}(s).
\]
Any minimax pair of policies is a Nash equilibrium.

\section*{5. Bellman Equations for Markov Games}

\paragraph{Value form.}
\[
		V^*(s)
		=
		\max_{p \in \Delta(A_1)}
		\min_{q \in \Delta(A_2)}
		\mathbb{E}_{a_1\sim p,\,a_2\sim q}\left[
				R(s,a_1,a_2)
				+
				\gamma P(s,a_1,a_2)V^*
		\right].
\]

\paragraph{Q-form.}
\[
		Q^*(s,a_1,a_2) = R(s,a_1,a_2) + \gamma P(s,a_1,a_2)V^*(s'),
\]
\[
		V^*(s)
		=
		\max_{p}\min_{q} \; p^\top Q^*(s) q.
\]

Thus, each state defines a \emph{matrix game} whose Nash value determines $V^*(s)$.

\section*{6. Dynamic Programming for Markov Games}

Finite-horizon backward induction:
\[
		V^*_k(s)
		=
		\max_p \min_q \mathbb{E}_{a_1,a_2}[R(s,a_1,a_2) + \gamma V^*_{k-1}(s')],
		\qquad V^*_0 = 0.
\]

\section*{7. Nash Q-Learning}

Hu \& Wellman (2003) generalize Q-learning to two-player zero-sum Markov games.

\subsection*{Algorithm}
\begin{itemize}
		\item Observe $s_t$, take action pair $a_t=(a_{1,t},a_{2,t})$.
		\item Receive reward $r_t$ and next state $s_{t+1}$.
		\item Compute the Nash value of the next-state matrix game:
    \[
						V_{t+1} = \max_p \min_q \; p^\top Q_t(s_{t+1}) q.
    \]
		\item Update:
    \[
						Q_{t+1}(s_t,a_t)
						\leftarrow
						(1-\alpha_t) Q_t(s_t,a_t)
						+ \alpha_t \left(r_t + \gamma V_{t+1}\right).
    \]
\end{itemize}

\subsection*{Convergence}
If $\sum_t \alpha_t = \infty$, $\sum_t \alpha_t^2 < \infty$, and all state–action pairs are visited
infinitely often, then:
\[
		Q_t \to Q^*.
\]

\section*{8. Algorithms for Finding Equilibria}

Common iterative schemes:
\begin{itemize}
		\item \textbf{Self-play:} each player best-responds to the other (may cycle).
		\item \textbf{Fictitious play:} each player best-responds to opponent’s empirical distribution. Converges
				asymptotically in zero-sum games.
		\item \textbf{Double oracle:} iteratively expand strategy sets with best responses.
		\item \textbf{Gradient / Mirror descent:} continuous updates for mixed strategies.
\end{itemize}

\paragraph{Challenges.}
\begin{itemize}
		\item Best responses may \emph{not} be Markovian even when optimal policies are.
		\item Reduction to normal-form games causes exponential blowup.
		\item MARL is generally non-stationary and violates single-agent assumptions.
\end{itemize}

\section*{9. Key Takeaways}
\begin{itemize}
		\item Multi-agent RL combines sequential decision-making with game theory.
		\item Nash equilibria, correlated equilibria, and minimax concepts are foundational.
		\item Markov games extend MDPs to competitive multi-agent settings with strong guarantees.
		\item Nash Q-learning generalizes value-based RL to adversarial environments.
		\item Many MARL algorithms face instability due to non-stationarity and equilibrium computation difficulty.
\end{itemize}
\end{document}
